% Generated by book/tools/notebook_to_tex.py. Do not edit by hand.

\chapter[작은 GPT 사전학습 (TinyStories CLM)]{작은 GPT 사전학습 (TinyStories Causal LM)}

\label{ch:24}

\markboth{작은 GPT 사전학습 (TinyStories CLM)}{작은 GPT 사전학습 (TinyStories CLM)}

\chaptermeta{24_gpt_tinystories/24_gpt_tinystories.ipynb}{https://colab.research.google.com/github/yoon-gu/neuqes-101/blob/master/24_gpt_tinystories/24_gpt_tinystories.ipynb}{작은 GPT2를 random init으로 만들고 직접 학습한 BPE 토크나이저와 TinyStories로 next-token 사전학습}



\index{1GPT@GPT}

\index{1decoder only@decoder-only}

\index{1causal language modeling@causal language modeling}

\index{1CausalLM@CausalLM}

\index{1next token prediction@next-token prediction}

\index{1CrossEntropyLoss@CrossEntropyLoss}

\index{1GPT2Config@GPT2Config}

\index{1GPT2LMHeadModel@GPT2LMHeadModel}

\index{1BPE@BPE}

\index{1ByteLevel@ByteLevel}

\index{1TinyStories@TinyStories}

\index{1DataCollatorForLanguageModeling@DataCollatorForLanguageModeling}

\index{1labels 100@labels=-100}

\index{1generate@generate}

\index{1temperature@temperature}

\index{1top k@top\_k}

\index{1top p@top\_p}

\index{0작은 GPT@작은 GPT}

\index{0디코더@디코더}

\index{0인과 언어 모델링@인과 언어 모델링}

\index{0다음 토큰 예측@다음 토큰 예측}

\index{0BPE 토크나이저@BPE 토크나이저}

\index{0직접 사전학습@직접 사전학습}

\index{0생성@생성}

\index{0샘플링@샘플링}

\index{1PreTrainedTokenizerFast@PreTrainedTokenizerFast}

\index{1BpeTrainer@BpeTrainer}

\index{1ByteLevelDecoder@ByteLevelDecoder}

\index{1GPT2LMHeadModel config@GPT2LMHeadModel(config)}

\index{1tie word embeddings@tie\_word\_embeddings}

\index{1causal attention@causal attention}

\index{1language modeling head@language modeling head}

\index{1random baseline@random baseline}

\index{1ln vocab@ln(vocab)}

\index{1perplexity@perplexity}

\index{1VRAM trace@VRAM trace}

\index{1model generate@model.generate}

\index{1sampling hyperparameter@sampling hyperparameter}

\index{1SFT@SFT}

\index{1alignment@alignment}

\index{1pretraining@pretraining}

\index{1response only loss@response-only loss}

\index{0TinyStories 사전학습@TinyStories 사전학습}

\index{0바이트 레벨 BPE@바이트 레벨 BPE}

\index{0언어 모델링 헤드@언어 모델링 헤드}

\index{0가중치 공유@가중치 공유}

\index{0무작위 초기화@무작위 초기화}

\index{0랜덤 기준선@랜덤 기준선}

\index{0학습 전 생성@학습 전 생성}

\index{0학습 후 생성@학습 후 생성}

\index{0응답 구간 손실@응답 구간 손실}



\section{변화 추적표}

\begin{booktable}{24장 변화 추적표 표}{tab:ch24-01}
\begin{adjustbox}{max width=\textwidth}
\begin{tabular}{@{}llllll@{}}
\toprule
Ch & 모델 & 토크나이저 & 데이터 & Output Head & Loss \\
\midrule

20 & 작은 BERT (영어, scratch) & \inlinecode{bert-base-uncased} (가져옴) & Wikitext-103 paragraphs & MLM head & \inlinecode{CrossEntropyLoss} (masked 15\%) \\
21 & \ref{ch:20}장 + 분류 헤드 & (\ref{ch:20}장과 동일) & Yelp 이진 (다른 도메인) & \inlinecode{Linear(H,\ 2)} & \inlinecode{CrossEntropyLoss} \\
22 & 작은 BERT (한국어, scratch) & \inlinecode{klue/bert-base} (가져옴) & 한국어 위키 paragraphs & MLM head & \inlinecode{CrossEntropyLoss} (masked 15\%) \\
23 & \ref{ch:22}장 + 분류 헤드 & (\ref{ch:22}장과 동일) & NSMC 이진 (다른 도메인) & \inlinecode{Linear(H,\ 2)} & \inlinecode{CrossEntropyLoss} \\
\textbf{24 \(\leftarrow\) 여기} & \textbf{작은 GPT2 (직접, scratch)} & \textbf{BPE (직접 학습, vocab=2048)} & \textbf{TinyStories 30K stories} & \textbf{\inlinecode{Linear(H,\ V)} (LM head, weight tied)} & \textbf{\inlinecode{CrossEntropyLoss} (next-token, 거의 모든 자리)} \\
25 (다음) & \inlinecode{gpt2} (124M, OpenAI WebText 사전학습) & BPE (GPT2 그대로) & TinyStories (\ref{ch:24}장과 동일) & \inlinecode{Linear(H,\ V)} (LM head 그대로) & \inlinecode{CrossEntropyLoss} (next-token) - \textbf{continual pretraining} \\
\bottomrule
\end{tabular}
\end{adjustbox}
\end{booktable}

전체 장 표는 \href{https://github.com/yoon-gu/neuqes-101\#챕터별-변화추적표}{루트 README} 를 참고하세요.

\begin{center}\rule{0.5\linewidth}{0.5pt}\end{center}

\section{\texorpdfstring{Phase 전환 --- Encoder (BERT) \(\to\) Decoder (GPT) 패러다임}{Phase 전환 --- Encoder (BERT) to Decoder (GPT) 패러다임}}

\ref{ch:07}--\ref{ch:23}장의 BERT 장들이 \emph{encoder + masked token 예측 + task head 부착 fine-tune} 패러다임이라면, Phase 4 (\ref{ch:24}--\ref{ch:31}장) 는 \emph{decoder + next-token 예측 + LM head 그대로 + SFT(behavior alignment)} 패러다임입니다. 본 장이 그 출발점입니다.

\begin{booktable}{24장 Phase 전환 - Encoder (BERT) to Decoder (GPT) 패러다임 표}{tab:ch24-02}
\begin{adjustbox}{max width=\textwidth}
\begin{tabular}{@{}lll@{}}
\toprule
축 & Phase 1·2·3 (BERT, \ref{ch:07}--\ref{ch:23}장) & \textbf{Phase 4 (GPT, \ref{ch:24}--\ref{ch:31}장)} \\
\midrule

본체 & Encoder (양방향 attention) & \textbf{Decoder (causal / masked attention)} \\
사전학습 task & MLM (가려진 토큰 예측) & \textbf{CausalLM (next-token 예측)} \\
학습 신호 위치 & 선택된 약 15\% 만 (\inlinecode{-100} 다수) & \textbf{거의 모든 토큰} (\inlinecode{-100} pad 만) \\
출력 head & task 별 부착 (\inlinecode{Linear(H,\ K)}) & \textbf{LM head (\inlinecode{Linear(H,\ V)}) 그대로} \\
Downstream 적응 & head 교체 + 본체 fine-tune (\emph{task 적응}) & \textbf{SFT (\emph{behavior alignment})} + alignment (DPO/GRPO) \\
“Fine-tune” 의미 & task 별 특화 & \textbf{prompt 만 바꿔도 다른 일} \\
\bottomrule
\end{tabular}
\end{adjustbox}
\end{booktable}

\begin{quote}
본 장은 그 \emph{출발점} --- 작은 GPT 를 처음부터 학습해 \emph{next-token 예측이 어떻게 generation 으로 이어지는지} 를 직접 봅니다. \ref{ch:25}장 (대규모 사전학습 \inlinecode{gpt2} 를 TinyStories 로 \textbf{continual pretraining}) / \ref{ch:28}장 (SFT) / \ref{ch:30}--\ref{ch:31}장 (DPO / GRPO) 가 같은 본체 위에 차곡차곡 쌓여 갑니다.
\end{quote}



\section{\texorpdfstring{Phase 3 \(\to\) Phase 4 다리 --- \emph{왜 갑자기 decoder 인가, 그리고 생성형 AI 의 등장}}{Phase 3 to Phase 4 다리 --- 왜 갑자기 decoder 인가, 그리고 생성형 AI 의 등장}}

\ref{ch:23}장 (한국어 BERT 분류) 에서 \ref{ch:24}장 (GPT) 로 넘어오면 \emph{왜 갑자기 encoder 를 버리고 decoder 로 가지?} 라는 의문이 듭니다. 그 사이를 잇는 큰 그림을 한 화면에 정리합니다.

\subsection{Transformer 의 세 갈래 --- encoder / decoder / encoder-decoder}

원래 Transformer (Vaswani et al. 2017, \emph{Attention Is All You Need}) 는 \emph{번역} 을 위한 \textbf{encoder-decoder} 구조였습니다. 이후 두 절반이 \emph{각자} 떨어져 나와 독립 계열이 됩니다.

\begin{booktable}{24장 Transformer 의 세 갈래 - encoder / decoder / encoder-decoder 표}{tab:ch24-03}
\begin{adjustbox}{max width=\textwidth}
\begin{tabular}{@{}lllll@{}}
\toprule
갈래 & attention & 대표 모델 & 잘하는 task & 본 커리큘럼 \\
\midrule

\textbf{Encoder-only} & 양방향 (bidirectional) & BERT, RoBERTa, DistilBERT, klue/bert-base & 분류·NER·추출형 QA (\emph{이해}) & \textbf{Phase 1-3 (\ref{ch:07}--\ref{ch:23}장)} \\
\textbf{Decoder-only} & causal (좌\(\to\)오 마스킹) & GPT-2/3/4, LLaMA, Mistral, KoGPT2 & generation·in-context learning·SFT/RLHF (\emph{생성}) & \textbf{Phase 4 (\ref{ch:24}--\ref{ch:31}장)} \\
\textbf{Encoder-Decoder (seq2seq)} & encoder 양방향 + decoder causal + \textbf{cross-attention} & T5, BART, mBART, KoBART, KE-T5 & 번역·요약·생성형 QA (\emph{변환}) & (본 커리큘럼은 미포함 --- 맥락만) \\
\bottomrule
\end{tabular}
\end{adjustbox}
\end{booktable}

\begin{quote}
본 커리큘럼은 \emph{encoder-only (이해)} \(\to\) \emph{decoder-only (생성)} 의 두 축을 직접 구현하며 잇습니다. \textbf{seq2seq} 는 \emph{입력 시퀀스를 출력 시퀀스로 변환} (번역·요약) 하는 제3의 길로, 두 절반을 다시 합치고 \emph{cross-attention} 으로 연결합니다 --- 본 커리큘럼에선 다루지 않지만 \emph{지형의 한 축} 으로 기억해 두세요.
\end{quote}

\subsection{왜 BERT 는 generation 이 어려운가}

BERT 의 \emph{양방향} attention 은 토큰 \(i\) 가 \emph{좌·우 전부} 를 봅니다. 그래서 \emph{다음 토큰을 좌\(\to\)오로 하나씩 뽑는} autoregressive 생성에는 부적합합니다 --- 미래 토큰을 이미 보고 있으니 “다음을 예측” 이 성립하지 않습니다. BERT 가 일부만 \texttt{{[}MASK{]}} 로 가려 복원하는 (MLM) 것도 이 \emph{양방향 cheating} 을 막기 위함이었습니다 (\ref{ch:20}장).

generation 을 하려면 둘 중 하나가 필요합니다:

\begin{itemize}
\tightlist
\item
  \textbf{causal masking} --- 미래를 가려 \emph{좌\(\to\)오 순차 생성} 을 가능케 함 \(\to\) \textbf{decoder (GPT), 본 장부터}
\item
  \textbf{반복 denoise} --- 양방향을 유지한 채 \emph{마스킹 비율을 일반화} 해 병렬 생성 \(\to\) \textbf{diffusion LM (후속 선택 주제)}
\end{itemize}

\begin{quote}
즉 BERT 의 양방향성은 \emph{이해} 엔 강점, \emph{순차 생성} 엔 약점입니다. 이 한 가지 차이가 decoder-only GPT와 diffusion LM이라는 두 생성 계열을 가릅니다. 본 원고에서는 먼저 GPT 계열을 끝까지 정리합니다.
\end{quote}

\subsection{attention 의 진화}

\begin{booktable}{24장 attention 의 진화 표}{tab:ch24-04}
\begin{adjustbox}{max width=\textwidth}
\begin{tabular}{@{}lll@{}}
\toprule
단계 & attention 방식 & 모델 \\
\midrule

양방향 self & 모든 위치가 서로를 봄 & BERT (encoder) \\
\textbf{causal masked self} & 과거(좌)만 봄 \(\to\) 미래 누출 차단 & \textbf{GPT (decoder), 본 장} \\
cross-attention & decoder 가 encoder 출력을 참조 & seq2seq (T5/BART) \\
\bottomrule
\end{tabular}
\end{adjustbox}
\end{booktable}

\subsection{생성형 AI 등장 타임라인 (큰 흐름)}

\begin{booktable}{24장 생성형 AI 등장 타임라인 (큰 흐름) 표}{tab:ch24-05}
\begin{adjustbox}{max width=\textwidth}
\begin{tabular}{@{}lll@{}}
\toprule
연도 & 사건 & 의미 \\
\midrule

2017 & Transformer (\emph{Attention Is All You Need}) & encoder-decoder, attention 의 출발 \\
2018 & BERT (encoder) / GPT-1 (decoder) & \emph{이해} 와 \emph{생성} 두 갈래 분기 \\
2019 & GPT-2 & 큰 decoder 의 generation 품질 --- 본 장의 reference 모델 \\
2020 & GPT-3 & \textbf{in-context learning} (예시만 줘도 task 수행, fine-tune 없이) \\
2022 & InstructGPT / ChatGPT & \textbf{SFT + RLHF} 로 \emph{지시를 따르는} 정렬 (\ref{ch:28}장, \ref{ch:30}장, \ref{ch:31}장의 주제) \\
2023+ & GPT-4 / LLaMA / Mistral / KoGPT 등 & decoder-only LLM 의 시대 \\
\bottomrule
\end{tabular}
\end{adjustbox}
\end{booktable}

\begin{quote}
본 커리큘럼 Phase 4 (\ref{ch:24}--\ref{ch:31}장) 가 이 타임라인을 \emph{압축 재현} 합니다 --- \emph{작은 GPT 사전학습 (본 장)} \(\to\) \emph{continual pretraining (\ref{ch:25}장)} \(\to\) \emph{SFT (\ref{ch:28}장)} \(\to\) \emph{alignment (\ref{ch:30}--\ref{ch:31}장)}. 2017-2022 의 흐름을 손으로 한 번 따라가 보는 셈입니다.
\end{quote}

\begin{quote}
\textbf{참고} --- \emph{T4 30분 룰 너머} 로 GPT/LLM scratch 학습을 더 키워 보고 싶다면 \href{https://github.com/FareedKhan-dev/train-llm-from-scratch}{FareedKhan-dev/train-llm-from-scratch} 가 좋은 출발점입니다. Transformer 를 PyTorch 로 직접 구현하고 \emph{13M \(\to\) 2B params} 까지 consumer GPU 로 학습하는 과정 + \textbf{GPU별 실용 모델 크기 표} (예: T4 16GB \(\approx\) 1.5-2B, RTX 4090 24GB \(\approx\) 4B, A100 40GB \(\approx\) 6-8B) 가 인상적입니다 --- 본 장의 약 3M 모델이 \emph{어디까지 커질 수 있는지} 의 감을 줍니다.
\end{quote}



\section{변경점: \ref{ch:23}장 대비}

\begin{booktable}{24장 변경점 요약}{tab:ch24-06}
\begin{adjustbox}{max width=\textwidth}
\begin{tabular}{@{}lll@{}}
\toprule
축 & \ref{ch:23}장 (한국어 BERT 분류 fine-tune) & \ref{ch:24}장 (GPT scratch + TinyStories) \\
\midrule

\textbf{모델 패밀리} & Encoder (\inlinecode{BertForSequenceClassification}) & \textbf{Decoder (\inlinecode{GPT2LMHeadModel})} \(\leftarrow\) \emph{Phase 전환의 핵심} \\
사전학습 task & MLM (\ref{ch:22}장 산출물 본체 + 분류 head) & \textbf{CausalLM (next-token, from scratch)} \\
토크나이저 & \inlinecode{klue/bert-base} WordPiece (가져옴, vocab 32K) & \textbf{BPE 직접 학습} (vocab 2,048) \\
데이터 & NSMC 한국어 영화 리뷰 (이진 라벨) & \textbf{TinyStories 영어 short stories} (라벨 없음) \\
Output head & \inlinecode{Linear(H,\ 2)} (새로 부착) & \textbf{\inlinecode{Linear(H,\ V)} LM head} (모델이 내장 + weight tied) \\
Loss & \inlinecode{CrossEntropyLoss} (분류, K=2) & \textbf{\inlinecode{CrossEntropyLoss} (next-token, K=V=2048)} \\
산출물 & 분류 정확도 & \textbf{generation 텍스트} (\inlinecode{model.generate()}) \\
\bottomrule
\end{tabular}
\end{adjustbox}
\end{booktable}

\begin{quote}
\textbf{변경점이 한꺼번에 많은 이유} --- Phase 가 바뀌는 \emph{전환 장} 라 \emph{축 자체} 가 새로 정의됩니다. \ref{ch:25}장 부터는 다시 \emph{한 가지 축} 만 바뀝니다 (\ref{ch:25}장: 본체 출발점 = scratch \(\to\) 사전학습 모델, 같은 \emph{continual pretraining} task / \ref{ch:26}장: 언어 (한국어 scratch) / \ref{ch:27}장: 한국어 continual pretraining / \ref{ch:28}장: 학습 단계 = pretraining \(\to\) SFT, \texttt{labels{[}prompt{]}\ =\ -100}).
\end{quote}



\section{\texorpdfstring{Loss --- \inlinecode{CrossEntropyLoss} (next-token)}{Loss --- CrossEntropyLoss (next-token)}}

수식은 MLM 의 CE 와 \emph{완전히 같음}. 다만 \emph{어느 자리에서 loss 가 계산되는가} 가 다릅니다.

\subsection{수식}

입력 토큰 시퀀스 \(x = (x_1, \dots, x_n)\) 에 대해, 각 위치 \(i\) 에서 \emph{그 다음 토큰} \(x_{i+1}\) 을 예측:

\begin{equation}
\label{eq:ch24-01}
L_{\text{CLM}} = -\frac{1}{n-1} \sum_{i=1}^{n-1} \log P(x_{i+1} \mid x_1, \dots, x_i)
\end{equation}

식~\eqref{eq:ch24-01}은 이 절에서 사용하는 핵심 관계를 정리한 것입니다.

\begin{itemize}
\tightlist
\item
  \(P(x_{i+1} \mid x_{\leq i})\): 모델이 \emph{지금까지 본 토큰만으로} 다음 토큰을 예측할 확률 (vocab 2,048 차원 softmax)
\item
  평균 분모 \(n-1\): pad 가 아닌 \emph{거의 모든} 자리에서 loss 계산 (MLM 의 15\% 와 대비)
\end{itemize}

\subsection{숫자로 감 잡기 (vocab=2048)}

\begin{booktable}{24장 숫자로 감 잡기 (vocab=2048) 표}{tab:ch24-07}
\begin{adjustbox}{max width=\textwidth}
\begin{tabular}{@{}lll@{}}
\toprule
모델 상태 & 정답 토큰 확률 & \(-\log p\) \\
\midrule

균등 추측 (random init 직후) & \(1/2048 \approx 4.88 \times 10^{-4}\) & \textbf{7.62} \(\leftarrow\) random baseline \\
약하게 학습 (정답 확률 0.02) & \(0.02\) & 3.91 \\
잘 학습된 작은 GPT (정답 확률 0.05-0.15) & \(0.05\) - \(0.15\) & \textbf{1.9 - 3.0} \(\leftarrow\) 이 장 목표 영역 \\
큰 사전학습 GPT (정답 확률 0.3+) & \(0.3\) & 1.20 \\
완벽 (정답 확률 1.0) & \(1.0\) & 0.00 \\
\bottomrule
\end{tabular}
\end{adjustbox}
\end{booktable}

\textbf{관전 포인트}:

\begin{itemize}
\tightlist
\item
  학습 첫 step loss 가 약 7.6 부근이면 random init 직후 \emph{균등 추측} 상태. 첫 100 step 안에 빠르게 떨어지면 vocab + 모델 정상.
\item
  목표는 \emph{vocab 후보를 좁히는} 단계 (약 2-3). TinyStories 의 단순한 어휘·문법 덕분에 3M 짜리 작은 모델로도 도달 가능.
\end{itemize}

\subsection{Perplexity}

\(\text{PPL} = e^{L}\) --- \emph{다음 토큰을 평균 몇 후보 중에서 고민하는가}:

\begin{booktable}{24장 Perplexity 표}{tab:ch24-08}
\begin{adjustbox}{max width=\textwidth}
\begin{tabular}{@{}lll@{}}
\toprule
CLM loss & PPL & 해석 \\
\midrule

7.62 & 2,048 & 균등 (전체 vocab) \\
4.0 & 55 & 약 50 개 후보 \\
2.5 & 12 & 약 12 개 후보 \(\leftarrow\) 본 장 목표 \\
1.0 & 2.7 & 거의 결정적 \\
\bottomrule
\end{tabular}
\end{adjustbox}
\end{booktable}

\begin{quote}
MLM 의 \(\ln(30522) \approx 10.33\) random baseline 과 같은 직관입니다. \emph{vocab 차원} 만 작아진 것 (2,048) 입니다.
\end{quote}



\section{\texorpdfstring{\inlinecode{labels\ =\ -100} thread 환기 --- MLM 의 \emph{15\% 만} vs CausalLM 의 \emph{거의 모든 자리}}{labels = -100 thread 환기 --- MLM 의 15\% 만 vs CausalLM 의 거의 모든 자리}}

\ref{ch:20}·\ref{ch:22}장의 MLM 에서 봤던 \inlinecode{labels\ =\ -100} ignore\_index 트릭이 GPT CausalLM 사전학습에서도 등장하지만, \textbf{적용 자리가 정반대} 입니다.

\begin{booktable}{24장 labels = -100 thread 환기 - MLM 의 15\% 만 vs CausalLM 의 거의 모든 자리 표}{tab:ch24-09}
\begin{adjustbox}{max width=\textwidth}
\begin{tabular}{@{}llll@{}}
\toprule
단계 & 장 & \inlinecode{labels} 구성 & loss 계산 자리 \\
\midrule

MLM 사전학습 & \ref{ch:20}장 (영어), \ref{ch:22}장 (한국어) & 선택된 약 15\% 만 원본 token id, 나머지 = \inlinecode{-100} & \emph{가려진 자리만} \\
\textbf{GPT CausalLM 사전학습} & \textbf{\ref{ch:24}장 (영어, 본 장), \ref{ch:26}장 (한국어)} & \textbf{\inlinecode{input\_ids.clone()} - pad 만 \inlinecode{-100}} & \textbf{거의 \emph{전 자리}} \\
SFT / Instruction Tuning & \ref{ch:28}장 (한국어 KoGPT2 SFT) & \textbf{prompt 부분 = \inlinecode{-100}}, \emph{답변 토큰만} 원본 id & \emph{답변 부분만} \\
\bottomrule
\end{tabular}
\end{adjustbox}
\end{booktable}

\begin{quote}
같은 \inlinecode{-100} 트릭, \emph{적용 자리만 정반대}. MLM 은 \emph{대부분을 가리고 일부만 학습}, GPT 사전학습은 \emph{거의 가리지 않음}, SFT 는 \emph{prompt 만 가림}. 한 step 에 학습되는 토큰 수만 봐도 \emph{GPT 사전학습은 MLM 대비 약 5-6배 효율} (15\% vs 거의 100\%).
\end{quote}

본 장에서는 \inlinecode{DataCollatorForLanguageModeling(mlm=False)} 이 자동으로 \inlinecode{labels\ =\ input\_ids.clone()} 을 만들어 줍니다 --- 뒤에 collator 출력 셀에서 직접 확인하겠습니다. \ref{ch:28}장의 \emph{왜 모델이 instruction 을 따라가게 되는가} 는 \emph{한 줄 코드 \texttt{labels{[}prompt\_mask{]}\ =\ -100}} 로 정확히 설명되는데, 그 코드를 이해할 토대가 \emph{이 장의 collator 출력} 입니다.



\section{\texorpdfstring{GPT 시대의 학습은 \emph{네 단계} --- 용어가 BERT 와 다릅니다}{GPT 시대의 학습은 네 단계 --- 용어가 BERT 와 다릅니다}}

\ref{ch:21}·\ref{ch:23}장에서 본 \emph{fine-tune} 은 \textbf{BERT 단계 의미} --- \emph{사전학습된 본체 + 새 task-specific head (\inlinecode{Linear(H,\ K)})}. 분류·회귀·QA 마다 다른 head, \emph{task 별 특화}. 한 모델 = 한 task.

Phase 4 GPT 시대는 \emph{fine-tune} 한 단어가 \emph{여러 의미} 로 섞여 쓰입니다. 학술적으로는 \textbf{네 단계} 로 분리됩니다.

\begin{booktable}{24장 GPT 시대의 학습은 네 단계 - 용어가 BERT 와 다릅니다 표}{tab:ch24-10}
\begin{adjustbox}{max width=\textwidth}
\begin{tabular}{@{}lllll@{}}
\toprule
단계 & 정확 용어 & 의미 & 학습 신호 & 본 커리큘럼 \\
\midrule

1 & \textbf{Pretraining} (사전학습) & 일반 코퍼스 위에 random init 본체부터 학습 & 모든 토큰 (\inlinecode{labels\ =\ input\_ids}) & \textbf{\ref{ch:24}장} (영어 scratch, TinyStories), \textbf{\ref{ch:26}장} (한국어 scratch) \\
2 & \textbf{Continual pretraining} (계속 사전학습 / continual learning) & \emph{사전학습된 본체} 를 \emph{새 데이터} 로 \emph{같은 CausalLM task} 더 학습. \textbf{head 그대로, task 그대로, 데이터만 새로} & 모든 토큰 (pretraining 과 동일) & \textbf{\ref{ch:25}장} (\inlinecode{gpt2} + TinyStories) \\
3 & \textbf{SFT} (Supervised Fine-Tuning / Instruction tuning) & instruction-response 쌍으로 \emph{행동 정렬}. \texttt{labels{[}prompt{]}\ =\ -100} 으로 답변 부분만 학습 & \textbf{답변 토큰만} & \textbf{\ref{ch:28}장} (KoGPT2 + KoAlpaca) \\
4 & \textbf{Alignment} (DPO / RLHF / GRPO) & preference 또는 verifier reward 로 \emph{선호 정렬} & preference log-likelihood ratio / RL advantage & \textbf{\ref{ch:30}장} (DPO), \textbf{\ref{ch:31}장} (GRPO) \\
\bottomrule
\end{tabular}
\end{adjustbox}
\end{booktable}

\textbf{세 가지 공통점} (모두 GPT 시대):

\begin{itemize}
\tightlist
\item
  \textbf{모델 클래스 그대로} --- \inlinecode{AutoModelForCausalLM} (BERT 처럼 task head 부착 안 함)
\item
  \textbf{출력 형식 그대로} --- 토큰 시퀀스
\item
  \textbf{학습 신호 종류 그대로} --- next-token CE (alignment 만 예외)
\end{itemize}

\textbf{다른 점은 \emph{데이터 형식} 과 \emph{어느 토큰에 학습 신호를 주는가}} --- \inlinecode{labels\ =\ -100} 자리만 변함.

\begin{booktable}{24장 GPT 시대의 학습은 네 단계 - 용어가 BERT 와 다릅니다 표}{tab:ch24-11}
\begin{adjustbox}{max width=\textwidth}
\begin{tabular}{@{}lll@{}}
\toprule
단계 & 데이터 & \inlinecode{labels\ =\ -100} 자리 \\
\midrule

Pretraining (\ref{ch:24}·\ref{ch:26}장) & 일반 텍스트 & pad 만 \\
\textbf{Continual pretraining (\ref{ch:25}장)} & \textbf{새 도메인 텍스트} & \textbf{pad 만 (Pretraining 과 동일)} \\
SFT (\ref{ch:28}장) & instruction + response & \textbf{prompt 부분} \\
Alignment (\ref{ch:30}·\ref{ch:31}장) & preference 쌍 / verifier reward & (RL 내부) \\
\bottomrule
\end{tabular}
\end{adjustbox}
\end{booktable}

\begin{quote}
\textbf{BERT 의 “fine-tune” 은 \emph{task 적응} 한 가지였지만, GPT 의 “fine-tune” 은 \emph{continual pretraining / SFT / alignment} 셋이 섞인 통칭}. 정확히 말하려면 단계별 용어를 구분합니다. 본 장은 단계 1 (사전학습) 그 자체. \ref{ch:25}장에서 단계 2 (continual pretraining) 가, \ref{ch:28}장에서 단계 3 (SFT --- 진짜 \emph{행동 정렬}) 이, \ref{ch:30}--\ref{ch:31}장에서 단계 4 (alignment) 가 본격 등장합니다.
\end{quote}

\begin{quote}
\emph{왜 GPT 모델 하나가 모든 task 를 해내는가} 의 답은 단계 3 (SFT) 부터 --- head 가 task 별로 분기하지 않으니 \emph{입력 프롬프트} 만 바꾸면 같은 모델이 다른 일을 합니다.
\end{quote}



\section{토크나이저 노트}

GPT-2 와 같은 종류 (byte-level BPE) 의 작은 vocab 을 직접 학습합니다. \ref{ch:19}장에서 봤던 WordPiece / WordLevel 토크나이저 학습 절차의 \emph{BPE 판}.

\begin{booktable}{24장 토크나이저 노트 표}{tab:ch24-12}
\begin{adjustbox}{max width=\textwidth}
\begin{tabular}{@{}lll@{}}
\toprule
토크나이저 & 학습 방식 & 등장 장 \\
\midrule

WordLevel & 공백 + 빈도 - 가장 단순 & \ref{ch:19}장 (직접 학습) \\
WordPiece & 빈도 기반 subword (BERT) & \ref{ch:07}--\ref{ch:23}장 (BERT 장들 - 가져옴), \ref{ch:19}장 (직접 학습) \\
\textbf{BPE (byte-level)} & \textbf{빈도 높은 byte 쌍 반복 병합 (GPT-2)} & \textbf{\ref{ch:24}장 (본 장, 직접 학습), \ref{ch:25}--\ref{ch:26}장 (GPT 장들)} \\
\bottomrule
\end{tabular}
\end{adjustbox}
\end{booktable}

\subsection{WordPiece vs BPE 의 결합 방식 차이}

같은 입력 \inlinecode{"unhappiness"} 에 대해:

\begin{itemize}
\tightlist
\item
  \textbf{WordPiece} (BERT): \texttt{{[}“un”,\ “\#\#happiness”{]}} 또는 \texttt{{[}“un”,\ “\#\#happy”,\ “\#\#ness”{]}} - 단어 \emph{중간} subword 에 \inlinecode{\#\#} 접두사. 단어 경계를 명시.
\item
  \textbf{BPE} (GPT-2): \texttt{{[}“un”,\ “happiness”{]}} 또는 \texttt{{[}“un”,\ “h”,\ “app”,\ “iness”{]}} - 접두사 없이 \emph{byte 시퀀스 그대로}. 단어 경계는 \emph{공백 자체가 한 byte} 로 처리.
\end{itemize}

byte-level BPE 의 핵심 장점: \emph{어떤 유니코드 문자열이든} (이모지, 한글, 특수 기호) UNK 없이 표현 가능 - 가장 작은 단위가 \emph{byte (256개)} 라 vocab 에 모든 byte 를 포함하면 \emph{완전 가역}.

\subsection{특수 토큰 컨벤션}

GPT-2 는 특수 토큰을 \emph{최소화} 합니다 - \texttt{\textless{}\textbar{}endoftext\textbar{}\textgreater{}} 하나만 사용 (bos = eos = pad 겸용). BERT 의 \texttt{{[}CLS{]}\ {[}SEP{]}\ {[}MASK{]}\ {[}PAD{]}\ {[}UNK{]}} 5종과 대비.

\begin{quote}
\ref{ch:19}장의 “토크나이저는 모델과 운명공동체” 원칙이 본 장에서도 유효 - vocab 2,048 의 BPE 를 직접 학습한 뒤, \emph{같은 vocab 으로 GPT 본체를 random init} 합니다. \ref{ch:25}장에서는 \emph{반대로} - \inlinecode{gpt2} (124M) 의 vocab 50,257 BPE 를 그대로 가져와 같은 TinyStories 데이터로 \textbf{continual pretraining}. 토크나이저 + 모델이 \emph{함께} 변하는 게 \ref{ch:24}--\ref{ch:25}장의 핵심 비교.
\end{quote}



\section{환경 준비}



\Needspace{5\baselineskip}
\begin{lstlisting}
%pip install -q -U transformers tokenizers datasets accelerate
\end{lstlisting}

\Needspace{11\baselineskip}
\begin{lstlisting}
import warnings
warnings.filterwarnings("ignore")

import math
import os
import random
import time

import numpy as np
import pandas as pd
import matplotlib.pyplot as plt
import torch

\end{lstlisting}

\noindent\emph{앞 블록은 필요한 라이브러리와 기본 설정을 준비하는 단계입니다. 이어지는 블록에서는 앞에서 만든 중간 값을 다음 계산으로 넘기는 단계로 넘어갑니다.}

\Needspace{11\baselineskip}
\begin{lstlisting}[firstnumber=14]
if torch.cuda.is_available():
    device = torch.device("cuda")
    device_name = torch.cuda.get_device_name(0)
    vram_gib = torch.cuda.get_device_properties(0).total_memory / 1024**3
    print(f"device     : cuda  ({device_name})")
    print(f"VRAM total : {vram_gib:.2f} GiB")
elif torch.backends.mps.is_available():
    device = torch.device("mps")
    print("device     : mps  (Apple Silicon)")
else:
    device = torch.device("cpu")
    print(
        "device     : cpu  (training will be very slow - "
        "Colab T4 recommended)"
    )

\end{lstlisting}

\noindent\emph{앞뒤 블록은 모두 앞에서 만든 중간 값을 다음 계산으로 넘기는 단계입니다. 길이를 나누어 같은 흐름을 단계별로 확인합니다.}

\Needspace{11\baselineskip}
\begin{lstlisting}[firstnumber=30]
print(f"torch      : {torch.__version__}")

SEED = 0
torch.manual_seed(SEED)
np.random.seed(SEED)
random.seed(SEED)

USE_FP16 = (device.type == "cuda")
print(f"use fp16   : {USE_FP16}")
\end{lstlisting}

\noindent\textbf{출력.}
\begin{bookoutputbox}
device     : cuda  (Tesla T4)
VRAM total : 14.56 GiB
torch      : 2.11.0+cu128
use fp16   : True
\end{bookoutputbox}
\par\vspace{0.9em}

\section{TinyStories 데이터 로드}

\inlinecode{roneneldan/TinyStories} 는 GPT-3.5 / GPT-4 가 \emph{4세 어린이가 이해할 단어만} 으로 생성한 짧은 영어 동화 약 2.1M 편 (Eldan \& Li 2023, arXiv:2305.07759). 어휘·문법이 단순해 \textbf{3-5M 파라미터} 짜리 작은 모델로도 grammatical 한 생성이 가능합니다.

학습 split 의 처음 \textbf{30,000 stories} 만 사용 (T4 30분 룰 안).



\Needspace{11\baselineskip}
\begin{lstlisting}
from datasets import load_dataset

N_TRAIN = 30_000
N_VAL   = 500

raw_train = load_dataset(
    "roneneldan/TinyStories",
    split=f"train[:{N_TRAIN}]",
)
raw_val   = load_dataset(
    "roneneldan/TinyStories",
    split=f"validation[:{N_VAL}]",
)
print("train:", raw_train)
print("val  :", raw_val)
print("\n=== sample story ===")
\end{lstlisting}

\noindent\emph{앞 블록은 필요한 라이브러리와 기본 설정을 준비하는 단계입니다. 이어지는 블록에서는 앞에서 만든 중간 값을 다음 계산으로 넘기는 단계로 넘어갑니다.}

\Needspace{5\baselineskip}
\begin{lstlisting}[firstnumber=17]
print(raw_train[0]["text"][:400])
\end{lstlisting}

\noindent\textbf{출력.}
\begin{bookoutputbox}
train: Dataset({
    features: ['text'],
    num_rows: 30000
})
val  : Dataset({
    features: ['text'],
    num_rows: 500
})

=== sample story ===
One day, a little girl named Lily found a needle in her room. She knew it w...

Lily went to her mom and said, "Mom, I found this needle. Can you share it...

To
\end{bookoutputbox}
\par\vspace{0.9em}

\section{BPE 토크나이저 직접 학습}

\inlinecode{tokenizers.BPE} + ByteLevel pre-tokenizer 로 vocab 2,048 의 BPE 를 코퍼스에서 직접 학습합니다. \ref{ch:19}장의 토크나이저 학습 절차와 같은 패턴 - 다른 점은 \emph{알고리즘} 만 (WordPiece \(\to\) BPE).



\Needspace{11\baselineskip}
\begin{lstlisting}
from tokenizers import Tokenizer
from tokenizers.models import BPE
from tokenizers.trainers import BpeTrainer
from tokenizers.pre_tokenizers import ByteLevel
from tokenizers.decoders import ByteLevel as ByteLevelDecoder
from transformers import PreTrainedTokenizerFast

VOCAB_SIZE = 2048
EOS = "<|endoftext|>"

\end{lstlisting}

\noindent\emph{앞 블록은 필요한 라이브러리와 기본 설정을 준비하는 단계입니다. 이어지는 블록에서는 텍스트를 모델 입력 텐서로 바꾸는 단계로 넘어갑니다.}

\Needspace{11\baselineskip}
\begin{lstlisting}[firstnumber=11]
bpe = Tokenizer(BPE(unk_token=None))
bpe.pre_tokenizer = ByteLevel(add_prefix_space=False)
bpe.decoder = ByteLevelDecoder()
trainer = BpeTrainer(
    vocab_size=VOCAB_SIZE,
    special_tokens=[EOS],
    initial_alphabet=ByteLevel.alphabet(),
    show_progress=True,
)

\end{lstlisting}

\noindent\emph{앞 블록은 텍스트를 모델 입력 텐서로 바꾸는 단계입니다. 이어지는 블록에서는 학습 설정을 만들고 학습 루프를 실행하는 단계로 넘어갑니다.}

\Needspace{11\baselineskip}
\begin{lstlisting}[firstnumber=21]
t0 = time.time()
bpe.train_from_iterator(
    (ex["text"] for ex in raw_train),
    trainer,
    length=len(raw_train),
)
print(
    f"BPE training done: {time.time()-t0:.1f}s, vocab="
    f"{bpe.get_vocab_size()}"
)

\end{lstlisting}

\noindent\emph{앞 블록은 학습 설정을 만들고 학습 루프를 실행하는 단계입니다. 이어지는 블록에서는 텍스트를 모델 입력 텐서로 바꾸는 단계로 넘어갑니다.}

\Needspace{11\baselineskip}
\begin{lstlisting}[firstnumber=32]
tokenizer = PreTrainedTokenizerFast(
    tokenizer_object=bpe,
    bos_token=EOS,
    eos_token=EOS,
    pad_token=EOS,
)

print("\n=== encode/decode demo ===")
sample = "Once upon a time, a little rabbit went to the forest."
enc = tokenizer(sample)
print(f"input      : {sample}")
print(f"ids        : {enc['input_ids']}")
print(
    f"tokens     : "
    f"{tokenizer.convert_ids_to_tokens(enc['input_ids'])}"
)
\end{lstlisting}

\noindent\emph{앞뒤 블록은 모두 텍스트를 모델 입력 텐서로 바꾸는 단계입니다. 길이를 나누어 같은 흐름을 단계별로 확인합니다.}

\Needspace{11\baselineskip}
\begin{lstlisting}[firstnumber=48]
print(
    f"decode     : "
    f"{tokenizer.decode(enc['input_ids'])}"
)
print(f"vocab_size : {tokenizer.vocab_size}")
print(
    f"eos_token  : {tokenizer.eos_token}  id="
    f"{tokenizer.eos_token_id}"
)
\end{lstlisting}

\begin{codeRead}
\inlinecode{tokenizers, transformers} 같은 주요 패키지를 불러옵니다.
\end{codeRead}

\noindent\textbf{출력.}
\begin{bookoutputbox}
BPE training done: 10.6s, vocab=2048

=== encode/decode demo ===
input      : Once upon a time, a little rabbit went to the forest.
ids        : [428, 440, 259, 394, 12, 259, 395, 1114, 464, 266, 263, 1081, 14]
tokens : ['Once', '<sp>upon', '<sp>a', '<sp>time', ',', '<sp>a', '<sp>littl...
decode     : Once upon a time, a little rabbit went to the forest.
vocab_size : 2048
eos_token  : <|endoftext|>  id=0
\end{bookoutputbox}
\par\vspace{0.9em}
\textbf{관전 포인트} --- \inlinecode{Once\ upon\ a\ time} 같이 TinyStories 에 \emph{자주 등장} 하는 표현은 \emph{적은 수의 토큰} 으로 압축, \inlinecode{rabbit} 처럼 덜 등장한 단어는 \emph{여러 byte 조각} 으로 분할되는 경향. vocab 2,048 은 작은 모델에 맞춘 \emph{최소한의 크기} 입니다.



\section{\texorpdfstring{토큰화 + \inlinecode{group\_texts} (HF 표준 CLM 전처리)}{토큰화 + group\_texts (HF 표준 CLM 전처리)}}

HuggingFace 의 causal LM 학습 표준 패턴 (\inlinecode{run\_clm.py}) 그대로:

\begin{enumerate}
\def\labelenumi{\arabic{enumi}.}
\tightlist
\item
  전체 코퍼스를 토큰화 (배치 단위)
\item
  각 story 끝에 \texttt{\textless{}\textbar{}endoftext\textbar{}\textgreater{}} 부착 (story 경계 표시)
\item
  모든 토큰을 이어붙여 1D 스트림으로 만든 뒤 \inlinecode{block\_size=128} 단위로 잘라 chunk 화
\item
  각 chunk 가 한 학습 sample - \inlinecode{DataCollatorForLanguageModeling(mlm=False)} 가 \inlinecode{labels\ =\ input\_ids} 를 자동으로 채워 next-token prediction loss 가 됨
\end{enumerate}

\ref{ch:20}·\ref{ch:22}장의 \inlinecode{group\_texts} 와 \emph{완전히 같은 패턴}. MLM 장들에선 \emph{마스킹} 만 추가됐다면, CausalLM 장에선 \emph{labels = input\_ids} 그대로.



\Needspace{11\baselineskip}
\begin{lstlisting}
BLOCK_SIZE = 128

def tokenize_fn(batch):
    return tokenizer(batch["text"])

tok_train = raw_train.map(
    tokenize_fn,
    batched=True,
    remove_columns=["text"],
    desc="tokenize train",
)
tok_val   = raw_val.map(
    tokenize_fn,
    batched=True,
    remove_columns=["text"],
    desc="tokenize val",
\end{lstlisting}

\noindent\emph{앞뒤 블록은 모두 텍스트를 모델 입력 텐서로 바꾸는 단계입니다. 길이를 나누어 같은 흐름을 단계별로 확인합니다.}

\Needspace{11\baselineskip}
\begin{lstlisting}[firstnumber=17]
)

def add_eos(batch):
    new_ids, new_mask = [], []
    for ids in batch["input_ids"]:
        ids = ids + [tokenizer.eos_token_id]
        new_ids.append(ids)
        new_mask.append([1] * len(ids))
    return {"input_ids": new_ids, "attention_mask": new_mask}

\end{lstlisting}

\noindent\emph{앞 블록은 텍스트를 모델 입력 텐서로 바꾸는 단계입니다. 이어지는 블록에서는 앞에서 만든 중간 값을 다음 계산으로 넘기는 단계로 넘어갑니다.}

\Needspace{11\baselineskip}
\begin{lstlisting}[firstnumber=27]
tok_train = tok_train.map(
    add_eos,
    batched=True,
    desc="add eos train",
)
tok_val   = tok_val.map(
    add_eos,
    batched=True,
    desc="add eos val",
)

\end{lstlisting}

\noindent\emph{앞뒤 블록은 모두 앞에서 만든 중간 값을 다음 계산으로 넘기는 단계입니다. 길이를 나누어 같은 흐름을 단계별로 확인합니다.}

\Needspace{10\baselineskip}
\begin{lstlisting}[firstnumber=38]
def group_texts(batch):
    concatenated = {k: sum(batch[k], []) for k in batch.keys()}
    total_len = len(concatenated["input_ids"])
    total_len = (total_len // BLOCK_SIZE) * BLOCK_SIZE
    return {
        k: [t[i : i + BLOCK_SIZE] for i in range(0, total_len, BLOCK_SIZE)]
        for k, t in concatenated.items()
    }

\end{lstlisting}

\noindent\emph{앞뒤 블록은 모두 앞에서 만든 중간 값을 다음 계산으로 넘기는 단계입니다. 길이를 나누어 같은 흐름을 단계별로 확인합니다.}

\Needspace{11\baselineskip}
\begin{lstlisting}[firstnumber=47]
lm_train = tok_train.map(
    group_texts,
    batched=True,
    desc="group train",
)
lm_val   = tok_val.map(
    group_texts,
    batched=True,
    desc="group val",
)

\end{lstlisting}

\noindent\emph{앞 블록은 앞에서 만든 중간 값을 다음 계산으로 넘기는 단계입니다. 이어지는 블록에서는 텍스트를 모델 입력 텐서로 바꾸는 단계로 넘어갑니다.}

\Needspace{11\baselineskip}
\begin{lstlisting}[firstnumber=58]
print(
    f"\ntrain chunks: {len(lm_train):,}  (block_size="
    f"{BLOCK_SIZE})"
)
print(f"val   chunks: {len(lm_val):,}")
print(
    f"approx. train tokens: "
    f"{len(lm_train) * BLOCK_SIZE / 1e6:.2f} M"
)
print("\nfirst chunk decode (first 200 chars):")
print(tokenizer.decode(lm_train[0]["input_ids"])[:200])
\end{lstlisting}

\noindent\textbf{출력.}
\begin{bookoutputbox}
train chunks: 57,973  (block_size=128)
val   chunks: 867
approx. train tokens: 7.42 M

first chunk decode (first 200 chars):
One day, a little girl named Lily found a needle in her room. She knew it w...
\end{bookoutputbox}
\par\vspace{0.9em}

\subsection{\texorpdfstring{Collator 가 만드는 \inlinecode{labels} 확인 - \emph{거의 모든 자리} 가 학습 신호}{Collator 가 만드는 labels 확인 - 거의 모든 자리 가 학습 신호}}

\inlinecode{DataCollatorForLanguageModeling(mlm=False)} 가 \emph{내부적으로} \inlinecode{labels\ =\ input\_ids.clone()} 을 만들어 \inlinecode{-100} 자리는 \emph{없거나 pad 토큰 자리만} 임을 직접 확인합니다. \ref{ch:20}·\ref{ch:22}장의 MLM collator 가 약 85\% 를 \inlinecode{-100} 으로 채웠던 것과 \emph{정확히 반대}.



\Needspace{11\baselineskip}
\begin{lstlisting}
from transformers import DataCollatorForLanguageModeling

collator_demo = DataCollatorForLanguageModeling(
    tokenizer=tokenizer,
    mlm=False,
)
demo_batch = collator_demo([lm_train[0], lm_train[1]])

input_ids = demo_batch["input_ids"]
labels = demo_batch["labels"]

print(f"input_ids shape: {tuple(input_ids.shape)}")
print(f"labels shape   : {tuple(labels.shape)}")

\end{lstlisting}

\noindent\emph{앞 블록은 필요한 라이브러리와 기본 설정을 준비하는 단계입니다. 이어지는 블록에서는 앞에서 만든 중간 값을 다음 계산으로 넘기는 단계로 넘어갑니다.}

\Needspace{11\baselineskip}
\begin{lstlisting}[firstnumber=15]
total = labels.numel()
n_ignored = (labels == -100).sum().item()
n_train_signal = total - n_ignored
print(
    f"\n=== 'labels = -100' thread - CausalLM vs MLM "
    f"comparison ==="
)
print(f"total positions      : {total}")
print(
    f"  ignored (-100)     : {n_ignored:>5d}  ("
    f"{100 * n_ignored / total:5.2f}%)"
)
print(
    f"  train signal       : {n_train_signal:>5d}  ("
    f"{100 * n_train_signal / total:5.2f}%)"
)
\end{lstlisting}

\noindent\emph{앞뒤 블록은 모두 앞에서 만든 중간 값을 다음 계산으로 넘기는 단계입니다. 길이를 나누어 같은 흐름을 단계별로 확인합니다.}

\Needspace{11\baselineskip}
\begin{lstlisting}[firstnumber=31]
print(
    f"\n[MLM (20장/22)]     approx. 85% = -100, 15% = "
    f"train signal"
)
print(
    f"[CausalLM (this ch)] {100 * n_ignored / total:5.2f}% = -100, "
    f"{100 * n_train_signal / total:5.2f}% = train signal  <- almost every position"
)
print(
    f"\n=> a single step's token-learning efficiency: GPT"
    f"pretrain is approx. 5-6x higher than MLM"
)

\end{lstlisting}

\noindent\emph{앞뒤 블록은 모두 앞에서 만든 중간 값을 다음 계산으로 넘기는 단계입니다. 길이를 나누어 같은 흐름을 단계별로 확인합니다.}

\Needspace{7\baselineskip}
\begin{lstlisting}[firstnumber=44]
identical = (input_ids == labels).sum().item()
print(
    f"\n(input_ids == labels) positions: {identical}/"
    f"{total}  - clone as-is"
)
\end{lstlisting}

\noindent\textbf{출력.}
\begin{bookoutputbox}
input_ids shape: (2, 128)
labels shape   : (2, 128)

=== 'labels = -100' thread - CausalLM vs MLM comparison ===
total positions      : 256
  ignored (-100)     :     1  ( 0.39%)
  train signal       :   255  (99.61%)

[MLM (Ch 20/22)]     approx. 85% = -100, 15% = train signal
[CausalLM (this ch)] 0.39% = -100, 99.61% = train signal <- almost every po...

=> a single step's token-learning efficiency: GPT pretrain is approx. 5-6x...

(input_ids == labels) positions: 255/256  - clone as-is
\end{bookoutputbox}
\par\vspace{0.9em}
\begin{quote}
\textbf{\inlinecode{-100} thread 환기} - MLM 은 \emph{마스킹된 자리만} 학습, CausalLM 은 \emph{거의 모든 자리} 학습. 같은 PyTorch \inlinecode{CrossEntropyLoss(ignore\_index=-100)} 트릭이 \emph{적용 자리만 정반대}. \ref{ch:28}장 (SFT) 에서는 \emph{prompt 자리만 -100} - 같은 트릭의 세 번째 적용. 그 한 줄 코드가 \emph{모델이 instruction 을 따라가게 만드는 핵심} 입니다.
\end{quote}

본 장의 collator 셋업이 \emph{그 토대} - \inlinecode{labels\ =\ input\_ids.clone()} 의 직관을 손에 익혀 두면 \ref{ch:28}장의 \texttt{labels{[}:prompt\_len{]}\ =\ -100} 가 단번에 이해됩니다.



\section{\texorpdfstring{\inlinecode{GPT2LMHeadModel} from scratch}{GPT2LMHeadModel from scratch}}

\inlinecode{GPT2Config} 의 핵심 필드만 작게 잡고 \emph{random init} (사전학습 X) 시작.

\begin{itemize}
\tightlist
\item
  \inlinecode{n\_layer=4,\ n\_head=4,\ n\_embd=256} \(\to\) 약 3M params, BERT 장들의 small DistilBERT 와 비슷한 스케일
\item
  \inlinecode{n\_positions\ =\ BLOCK\_SIZE\ =\ 128} - 학습한 만큼만 context 사용
\item
  bos / eos / pad token id 를 토크나이저와 동기화
\item
  \inlinecode{tie\_word\_embeddings=True} (기본) - LM head 와 input embedding 의 weight 를 공유 \(\to\) 파라미터 절약
\end{itemize}

\subsection{BERT 와의 차이가 코드로 드러나는 곳}

\begin{itemize}
\tightlist
\item
  \inlinecode{BertForMaskedLM} 이 아니라 \inlinecode{GPT2LMHeadModel} - 클래스 자체가 \emph{causal attention 내장}
\item
  \inlinecode{from\_pretrained(...)} 없이 \inlinecode{GPT2LMHeadModel(config)} - 무작위 초기화 from scratch (\ref{ch:20}·\ref{ch:22}장의 \inlinecode{BertForMaskedLM(config)} 와 같은 패턴, \emph{모델 패밀리만} 다름)
\end{itemize}



\Needspace{11\baselineskip}
\begin{lstlisting}
from transformers import GPT2Config, GPT2LMHeadModel

config = GPT2Config(
    vocab_size=tokenizer.vocab_size,
    n_positions=BLOCK_SIZE,
    n_embd=256,
    n_layer=4,
    n_head=4,
    bos_token_id=tokenizer.bos_token_id,
    eos_token_id=tokenizer.eos_token_id,
    pad_token_id=tokenizer.pad_token_id,
    activation_function="gelu_new",
    resid_pdrop=0.1, embd_pdrop=0.1, attn_pdrop=0.1,
)

\end{lstlisting}

\noindent\emph{앞 블록은 필요한 라이브러리와 기본 설정을 준비하는 단계입니다. 이어지는 블록에서는 앞에서 만든 중간 값을 다음 계산으로 넘기는 단계로 넘어갑니다.}

\Needspace{11\baselineskip}
\begin{lstlisting}[firstnumber=16]
model = GPT2LMHeadModel(config).to(device)
n_params = model.num_parameters()
print(f"#params           : {n_params/1e6:.2f} M")
print(
    f"weight tying      : "
    f"{config.tie_word_embeddings}  (lm_head <-> wte shared)"
)
print(
    f"fp32 weight size  : "
    f"{n_params * 4 / 1024**2:.2f} MiB"
)
print(f"\nmodel: {type(model).__name__}")
print(
    f"  - body : "
    f"{type(model.transformer).__name__}  (Decoder, causal attention)"
)
\end{lstlisting}

\noindent\emph{앞뒤 블록은 모두 앞에서 만든 중간 값을 다음 계산으로 넘기는 단계입니다. 길이를 나누어 같은 흐름을 단계별로 확인합니다.}

\Needspace{6\baselineskip}
\begin{lstlisting}[firstnumber=32]
print(
    f"  - head : {type(model.lm_head).__name__}(in={model.lm_head.in_features}, out="
    f"{model.lm_head.out_features})"
)
\end{lstlisting}

\noindent\textbf{출력.}
\begin{bookoutputbox}
#params           : 3.72 M
weight tying      : True  (lm_head <-> wte shared)
fp32 weight size  : 14.18 MiB

model: GPT2LMHeadModel
  - body : GPT2Model  (Decoder, causal attention)
  - head : Linear(in=256, out=2048)
\end{bookoutputbox}
\par\vspace{0.9em}

\section{\texorpdfstring{학습 \emph{전} generation - 비교 기준선 (random init baseline)}{학습 전 generation - 비교 기준선 (random init baseline)}}

\ref{ch:20}·\ref{ch:22}장의 \emph{사전학습 전 {[}MASK{]} top-5 후보} 와 같은 역할. random init 모델은 통계적으로 \emph{어느 토큰이든 거의 균등한 확률} 로 뽑으니, 생성 텍스트가 \emph{영어와 거리가 먼 byte 조각 / 의미 없는 짧은 단어 나열} 이 나옵니다.

같은 prompt 와 sampling 설정을 학습 \emph{전 / 후} 모두에서 호출 \(\to\) loss 곡선 없이도 \emph{학습이 본체에 무엇을 새겼는가} 가 한 화면에 드러납니다.



\Needspace{11\baselineskip}
\begin{lstlisting}
PROMPTS = [
    "Once upon a time,",
    "The little girl",
    "A big dog",
]
GEN_KWARGS = dict(
    max_new_tokens=60,
    do_sample=True,
    temperature=0.8,
    top_k=50,
)

\end{lstlisting}

\noindent\emph{앞 블록은 앞에서 만든 중간 값을 다음 계산으로 넘기는 단계입니다. 이어지는 블록에서는 텍스트를 모델 입력 텐서로 바꾸는 단계로 넘어갑니다.}

\Needspace{11\baselineskip}
\begin{lstlisting}[firstnumber=13]
@torch.no_grad()
def generate_text(active_model, prompt: str, gen_tokenizer=None, **kwargs):
    tok = gen_tokenizer if gen_tokenizer is not None else tokenizer
    enc = tok(prompt, return_tensors="pt").to(active_model.device)
    out = active_model.generate(
        **enc,
        pad_token_id=tok.pad_token_id,
        eos_token_id=tok.eos_token_id,
        **kwargs,
    )
    return tok.decode(out[0], skip_special_tokens=True)

\end{lstlisting}

\noindent\emph{앞 블록은 텍스트를 모델 입력 텐서로 바꾸는 단계입니다. 이어지는 블록에서는 앞에서 만든 중간 값을 다음 계산으로 넘기는 단계로 넘어갑니다.}

\Needspace{11\baselineskip}
\begin{lstlisting}[firstnumber=25]
torch.manual_seed(SEED)
model.eval()
before_outputs = []
print("=" * 70)
print(
    "UNTRAINED model - generation from random initial "
    "weights"
)
print("=" * 70)
for p in PROMPTS:
    text = generate_text(model, p, **GEN_KWARGS)
    before_outputs.append(text)
    print(f"\n[prompt] {p}")
    print(text)
\end{lstlisting}

\noindent\textbf{출력.}
\begin{bookoutputbox}
======================================================================
UNTRAINED model - generation from random initial weights
======================================================================

[prompt] Once upon a time,
Once upon a time,ushinkush min is wondered5 cruallyked bed farmer smo wonde...

[prompt] The little girl
The little girlakak everyush Sarahgged:un't different different# gl keepner...

[prompt] A big dog
A big dog cle music hisftere learnedpe fam pullve bat batinin paper paper t...
\end{bookoutputbox}
\par\vspace{0.9em}
\textbf{관전 포인트} - 학습 전 출력은 \emph{무작위 토큰 나열} (반복되는 짧은 byte 조각, 의미 없는 단어들). \ref{ch:20}·\ref{ch:22}장의 \emph{학습 전 {[}MASK{]} top-5} 가 \emph{the / a / of / , / .} 같은 통계적 빈도 토큰이었던 것과 같은 현상의 \emph{generation 판} 입니다. 학습 후 출력과 \emph{나란히 비교} 하면 사전학습이 본체에 \emph{next-token 분포} 를 새긴 증거를 직접 보게 됩니다.



\section{\texorpdfstring{\inlinecode{Trainer} 로 사전학습}{Trainer 로 사전학습}}

BERT 장들 (\ref{ch:20}·\ref{ch:22}장) 과 \emph{완전히 같은} Trainer 패턴 - 바뀌는 건 모델 클래스와 collator 의 \inlinecode{mlm=False} 두 곳.

\begin{itemize}
\tightlist
\item
  \inlinecode{DataCollatorForLanguageModeling(mlm=False)} \(\to\) \inlinecode{labels\ =\ input\_ids} (next-token prediction)
\item
  \inlinecode{max\_steps=1500}, \inlinecode{batch\_size=32}, \inlinecode{fp16=True} - T4 약 1분
\item
  \inlinecode{eval\_steps=150} 으로 train / val loss 추이 관찰
\end{itemize}



\Needspace{11\baselineskip}
\begin{lstlisting}
from transformers import (DataCollatorForLanguageModeling, Trainer,
                          TrainingArguments, TrainerCallback)

collator = DataCollatorForLanguageModeling(
    tokenizer=tokenizer,
    mlm=False,
)

args = TrainingArguments(
    output_dir="./out_gpt_tinystories",
    max_steps=1500,
    per_device_train_batch_size=32,
    per_device_eval_batch_size=32,
    learning_rate=3e-4,
    weight_decay=0.1,
    adam_beta1=0.9, adam_beta2=0.95,
\end{lstlisting}

\noindent\emph{앞 블록은 필요한 라이브러리와 기본 설정을 준비하는 단계입니다. 이어지는 블록에서는 앞에서 만든 중간 값을 다음 계산으로 넘기는 단계로 넘어갑니다.}

\Needspace{11\baselineskip}
\begin{lstlisting}[firstnumber=17]
    warmup_steps=100,
    lr_scheduler_type="cosine",
    max_grad_norm=1.0,
    fp16=USE_FP16,
    logging_steps=50,
    eval_strategy="steps",
    eval_steps=150,
    save_strategy="no",
    report_to="none",
    dataloader_num_workers=2,
    dataloader_pin_memory=True,
    seed=SEED,
)

\end{lstlisting}

\noindent\emph{앞 블록은 앞에서 만든 중간 값을 다음 계산으로 넘기는 단계입니다. 이어지는 블록에서는 학습 설정을 만들고 학습 루프를 실행하는 단계로 넘어갑니다.}

\Needspace{11\baselineskip}
\begin{lstlisting}[firstnumber=31]
class VRAMCallback(TrainerCallback):
    '''step 별 peak VRAM 기록 (로깅 윈도우 단위로 reset). CUDA 에서만 유효.'''

    def __init__(self):
        self.steps, self.peak_MiB = [], []

    def on_train_begin(self, args, state, control, **kwargs):
        if torch.cuda.is_available():
            torch.cuda.reset_peak_memory_stats()

\end{lstlisting}

\noindent\emph{앞 블록은 학습 설정을 만들고 학습 루프를 실행하는 단계입니다. 이어지는 블록에서는 앞에서 만든 중간 값을 다음 계산으로 넘기는 단계로 넘어갑니다.}

\Needspace{10\baselineskip}
\begin{lstlisting}[firstnumber=41]
    def on_log(self, args, state, control, logs=None, **kwargs):
        if torch.cuda.is_available():
            peak = torch.cuda.max_memory_allocated() / 1024**2
            self.steps.append(state.global_step)
            self.peak_MiB.append(peak)
            torch.cuda.reset_peak_memory_stats()

vram_cb = VRAMCallback()

\end{lstlisting}

\noindent\emph{앞 블록은 앞에서 만든 중간 값을 다음 계산으로 넘기는 단계입니다. 이어지는 블록에서는 텍스트를 모델 입력 텐서로 바꾸는 단계로 넘어갑니다.}

\Needspace{11\baselineskip}
\begin{lstlisting}[firstnumber=50]
trainer = Trainer(
    model=model,
    args=args,
    train_dataset=lm_train,
    eval_dataset=lm_val,
    data_collator=collator,
    callbacks=[vram_cb],
)

t0 = time.time()
train_out = trainer.train()
elapsed = time.time() - t0

\end{lstlisting}

\noindent\emph{앞뒤 블록은 모두 텍스트를 모델 입력 텐서로 바꾸는 단계입니다. 길이를 나누어 같은 흐름을 단계별로 확인합니다.}

\Needspace{11\baselineskip}
\begin{lstlisting}[firstnumber=63]
print(f"\n=== training summary ===")
print(f"elapsed       : {elapsed/60:.2f} min")
print(f"global_step   : {train_out.global_step}")
print(f"train_loss    : {train_out.training_loss:.4f}")
print(
    f"random baseline (ln vocab): "
    f"{math.log(tokenizer.vocab_size):.4f}"
)
if torch.cuda.is_available():
    print(
        f"final peak    : "
        f"{torch.cuda.max_memory_allocated()/1024**2:.0f} MiB"
    )
\end{lstlisting}

\noindent\textbf{출력.}
\begin{bookoutputbox}
[transformers] `loss_type=None` was set in the config but it is unrecognize...
Step  Training Loss  Validation Loss
----  -------------  ---------------
150   4.943553       4.512319
300   4.164156       3.928088
450   3.859346       3.635721
600   3.661513       3.438943
750   3.525258       3.297616
900   3.410508       3.207347
1050  3.381458       3.144733
1200  3.332278       3.108168
1350  3.313862       3.091538
1500  3.305599       3.088617

=== training summary ===
elapsed       : 0.87 min
...
\end{bookoutputbox}
\par\vspace{0.9em}
\Needspace{11\baselineskip}
\begin{lstlisting}
# loss curve + VRAM trace
log = trainer.state.log_history
train_pts = [(r["step"], r["loss"]) for r in log if "loss" in r and "eval_loss" not in r]
eval_pts  = [(r["step"], r["eval_loss"]) for r in log if "eval_loss" in r]

fig, (ax1, ax2) = plt.subplots(1, 2, figsize=(13, 4))

# loss
ax1.plot([s for s, _ in train_pts], [l for _, l in train_pts], "-",
         color="tab:blue", alpha=0.6, label="train")
\end{lstlisting}

\noindent\emph{앞 블록은 학습 설정을 만들고 학습 루프를 실행하는 단계입니다. 이어지는 블록에서는 텍스트를 모델 입력 텐서로 바꾸는 단계로 넘어갑니다.}

\Needspace{10\baselineskip}
\begin{lstlisting}[firstnumber=11]
if eval_pts:
    ax1.plot([s for s, _ in eval_pts], [l for _, l in eval_pts], "s-",
             color="tab:red", label="eval")
ax1.axhline(math.log(tokenizer.vocab_size), ls=":", color="gray",
            label=f"uniform baseline = ln({tokenizer.vocab_size}) approx. {math.log(tokenizer.vocab_size):.2f}")
ax1.set_xlabel("step"); ax1.set_ylabel("cross-entropy loss")
ax1.set_title("TinyGPT-2 on TinyStories - loss")
ax1.grid(True, alpha=0.3); ax1.legend()

\end{lstlisting}

\noindent\emph{앞 블록은 텍스트를 모델 입력 텐서로 바꾸는 단계입니다. 이어지는 블록에서는 숫자 결과를 그림으로 확인하는 단계로 넘어갑니다.}

\Needspace{11\baselineskip}
\begin{lstlisting}[firstnumber=20]
if vram_cb.steps:
    ax2.plot(vram_cb.steps, vram_cb.peak_MiB, "o-", color="tab:green",
             label="peak VRAM (per log window)")
    ax2.set_title(f"VRAM trace  (bs=32, fp16, n_pos={BLOCK_SIZE})")
else:
    ax2.text(0.5, 0.5, "VRAM trace available on CUDA only",
             ha="center", va="center", transform=ax2.transAxes)
    ax2.set_title("VRAM trace - CUDA only")
ax2.set_xlabel("step"); ax2.set_ylabel("VRAM (MiB)")
ax2.grid(True, alpha=0.3); ax2.legend()

plt.tight_layout(); plt.show()
\end{lstlisting}

\noindent\textbf{출력.}
\begin{bookoutputbox}
<Figure size 1300x400 with 2 Axes>
\end{bookoutputbox}
\par\vspace{0.9em}
\textbf{관전 포인트} - 학습 첫 step loss 가 약 7.6 (random baseline \inlinecode{ln(2048)}) 부근에서 시작해 \emph{수백 step 안에 약 4-5} 로 빠르게 떨어지고, 1500 step 끝에 \emph{train\_loss 약 3.8} 부근이면 정상. perplexity 로 환산하면 vocab 2,048 중 \emph{약 12-20 개 후보} 로 좁힌 상태 - 다음 토큰을 \emph{어느 정도 결정적} 으로 뽑는 수준.



\section{\texorpdfstring{학습 \emph{후} generation + before/after 비교}{학습 후 generation + before/after 비교}}

같은 \inlinecode{PROMPTS\ /\ GEN\_KWARGS} 로 학습 후 모델에서 다시 생성하고, §5 의 학습 전 결과와 나란히 비교합니다. \textbf{이 장의 합격 기준}: 학습 후 텍스트가 \emph{전} 보다 명확히 \emph{영어 문장} 에 가까워졌는가 - \ref{ch:20}·\ref{ch:22}장의 \emph{사전·사후 {[}MASK{]} top-5} 비교의 \emph{generation 판}.



\Needspace{11\baselineskip}
\begin{lstlisting}
torch.manual_seed(SEED)
model.eval()
after_outputs = []
print("=" * 70)
print("TRAINED model - generation after Trainer.train()")
print("=" * 70)
for p in PROMPTS:
    text = generate_text(model, p, **GEN_KWARGS)
    after_outputs.append(text)
    print(f"\n[prompt] {p}")
    print(text)
\end{lstlisting}

\noindent\textbf{출력.}
\begin{bookoutputbox}
======================================================================
TRAINED model - generation after Trainer.train()
======================================================================

[prompt] Once upon a time,
Once upon a time, there was a girl named Lily. She loved to play with her f...

One day, Lily saw a small house, a boy named Timmy. He was so happy because...

[prompt] The little girl
The little girl had been a wonderful time. It was so happy to see the park....

[prompt] A big dog
A big dog, but they could go in the park. They ran away and the truck. The...

"But we can't get my mouth. I can play with you."

"It's okay, it is not very curious. He is not
\end{bookoutputbox}
\par\vspace{0.9em}
\Needspace{11\baselineskip}
\begin{lstlisting}
print("=" * 78)
print(
    "BEFORE (random init) vs AFTER (trained on "
    "TinyStories 30K)"
)
print("=" * 78)
for p, before, after in zip(PROMPTS, before_outputs, after_outputs):
    print(f"\nPROMPT  : {p}")
    print("-" * 78)
    print(f"BEFORE  : {before[len(p):].strip()[:280]}")
    print(f"AFTER   : {after[len(p):].strip()[:280]}")
\end{lstlisting}

\noindent\textbf{출력.}
\begin{bookoutputbox}
==============================================================================
BEFORE (random init) vs AFTER (trained on TinyStories 30K)
==============================================================================

PROMPT  : Once upon a time,
------------------------------------------------------------------------------
BEFORE : ushinkush min is wondered5 cruallyked bed farmer smo wonder smo dr...
AFTER : there was a girl named Lily. She loved to play with her friends. Sh...

One day, Lily saw a small house, a boy named Timmy. He was so happy because...

PROMPT  : The little girl
------------------------------------------------------------------------------
BEFORE : akak everyush Sarahgged:un't different different# gl keepner Graie...
AFTER : had been a wonderful time. It was so happy to see the park. She tha...

...
\end{bookoutputbox}
\par\vspace{0.9em}
\textbf{해석 가이드 - 사전학습이 만든 차이}

\begin{itemize}
\tightlist
\item
  \textbf{BEFORE (random init)}: \emph{영어와 거리가 먼 byte 조각 / 의미 없는 짧은 단어 반복}. logits 가 random 초기값이라 sampling 이 통계적 빈도 토큰들 사이에서만 흔들림.
\item
  \textbf{AFTER (TinyStories 30K \(\times\) 1500 steps)}: \emph{말이 되는 영어 문장} - 짧지만 \emph{주어 + 동사 + 목적어} 구조, \emph{동화 풍 어휘} (rabbit, forest, friend, mom, happy, ...). 완벽하진 않아도 \emph{학습이 본체에 next-token 분포를 새긴 증거} 가 한 줄에서 명확.
\end{itemize}

\begin{quote}
\ref{ch:20}·\ref{ch:22}장의 \emph{사전·사후 {[}MASK{]} top-5} 비교에서 \texttt{{[}MASK{]}} 자리에 \emph{the / a / of} 같은 빈도 토큰만 뽑히던 random init 모델이, 학습 후엔 \emph{문맥에 맞는 정답 토큰} 을 top-5 에 담아내던 그 변화의 \emph{generation 판} 입니다.
\end{quote}



\subsection{\texorpdfstring{Reference 비교 - \inlinecode{gpt2} (124M, OpenAI WebText) 의 같은 prompt generation}{Reference 비교 - gpt2 (124M, OpenAI WebText) 의 같은 prompt generation}}

같은 prompt 3개를 \emph{학습이 충분히 잘 된} 표준 \inlinecode{gpt2} (124M params, WebText 약 40GB 사전학습) 에 넣어 \emph{우리 작은 GPT (약 3M, TinyStories 30K)} 와 격차를 직접 비교. \ref{ch:20}장의 \emph{3-way {[}MASK{]} top-5 비교} (before / ours / \inlinecode{bert-base-uncased}) 와 같은 패턴.

T4 에서 약 1분 추가. 데이터·파라미터 격차가 generation 품질의 격차로 어떻게 드러나는지 한 화면에.



\Needspace{11\baselineskip}
\begin{lstlisting}
from transformers import AutoTokenizer, AutoModelForCausalLM

print(
    "loading reference gpt2 (124M, OpenAI WebText "
    "pretraining)..."
)
ref_tok = AutoTokenizer.from_pretrained("gpt2")
ref_tok.pad_token = ref_tok.eos_token
ref_model = AutoModelForCausalLM.from_pretrained("gpt2").to(device).eval()
print(f"  vocab_size : {ref_tok.vocab_size:,}")
print(
    f"  #params    : "
    f"{ref_model.num_parameters()/1e6:.1f} M"
)

\end{lstlisting}

\noindent\emph{앞 블록은 필요한 라이브러리와 기본 설정을 준비하는 단계입니다. 이어지는 블록에서는 텍스트를 모델 입력 텐서로 바꾸는 단계로 넘어갑니다.}

\Needspace{11\baselineskip}
\begin{lstlisting}[firstnumber=16]
torch.manual_seed(SEED)
ref_outputs = []
print("\n" + "=" * 70)
print(
    "REFERENCE gpt2 (124M, WebText) - generation on same "
    "prompts"
)
print("=" * 70)
for p in PROMPTS:
    text = generate_text(
        ref_model,
        p,
        gen_tokenizer=ref_tok,
        **GEN_KWARGS,
    )
    ref_outputs.append(text)
\end{lstlisting}

\noindent\emph{앞 블록은 텍스트를 모델 입력 텐서로 바꾸는 단계입니다. 이어지는 블록에서는 앞에서 만든 중간 값을 다음 계산으로 넘기는 단계로 넘어갑니다.}

\Needspace{8\baselineskip}
\begin{lstlisting}[firstnumber=32]
    print(f"\n[prompt] {p}")
    print(text)

del ref_model
if torch.cuda.is_available():
    torch.cuda.empty_cache()
\end{lstlisting}

\noindent\textbf{출력.}
\begin{bookoutputbox}
loading reference gpt2 (124M, OpenAI WebText pretraining)...
  vocab_size : 50,257
  #params    : 124.4 M

======================================================================
REFERENCE gpt2 (124M, WebText) - generation on same prompts
======================================================================

[prompt] Once upon a time,
Once upon a time, if you don't know what your country's government is doing...

In the last few months, I've traveled to dozens of countries around the wor...

My new book — the Making of a Better World Order:

[prompt] The little girl
...
\end{bookoutputbox}
\par\vspace{0.9em}
\Needspace{11\baselineskip}
\begin{lstlisting}
print("=" * 78)
print(
    "3-way comparison: BEFORE (random) / OURS (3M, "
    "TinyStories 30K) / REF (gpt2 124M, WebText)"
)
print("=" * 78)
for p, before, after, ref in zip(PROMPTS, before_outputs, after_outputs, ref_outputs):
    print(f"\nPROMPT : {p}")
    print("-" * 78)
    print(f"BEFORE : {before[len(p):].strip()[:240]}")
    print(f"OURS   : {after[len(p):].strip()[:240]}")
    print(f"REF    : {ref[len(p):].strip()[:240]}")
\end{lstlisting}

\noindent\textbf{출력.}
\begin{bookoutputbox}
==============================================================================
3-way comparison: BEFORE (random) / OURS (3M, TinyStories 30K) / REF (gpt2...
==============================================================================

PROMPT : Once upon a time,
------------------------------------------------------------------------------
BEFORE : ushinkush min is wondered5 cruallyked bed farmer smo wonder smo dr...
OURS : there was a girl named Lily. She loved to play with her friends. She...

One day, Lily saw a small house, a boy named Timmy. He was so happy because...
REF : if you don't know what your country's government is doing, you can fi...

In the last few months, I've traveled to dozens of countries around the wor...

My new book — the Making of a Better World Orde

...
\end{bookoutputbox}
\par\vspace{0.9em}
\textbf{해석 가이드 - 데이터·파라미터 규모가 만든 격차}

\begin{itemize}
\tightlist
\item
  \textbf{BEFORE (random)}: 영어와 거리 먼 byte 조각.
\item
  \textbf{OURS (3M, TinyStories 30K \(\times\) 1500 steps)}: \emph{동화 풍 단순 영어} - 어휘는 동화 도메인에 강하지만 (rabbit, forest, mom, friend, ...) \emph{복잡한 문장 구조 / 추상적 어휘} 는 약함.
\item
  \textbf{REF (gpt2 124M, WebText 약 40GB)}: \emph{다양한 도메인 어휘 + 자연스러운 문장 흐름} - 같은 prompt 에 대해 \emph{동화풍이 아닌 일반 산문 / 뉴스 / 대화} 등 다양한 톤. 학습 데이터 분포 (WebText) 의 다양성이 generation 다양성으로 직결.
\end{itemize}

\begin{quote}
\textbf{세 모델의 격차가 정확히 \emph{모델 크기 + 데이터 크기 + 데이터 다양성} 의 격차} - 우리 작은 GPT (3M, TinyStories 30K stories) \(\to\) reference \inlinecode{gpt2} (124M, WebText 약 40GB) 사이에 \emph{파라미터 약 40배, 데이터 규모 약 수천 배, 도메인 다양성 격차}. 그게 generation 의 \emph{질적 차이} 로 정확히 드러납니다.
\end{quote}

\begin{quote}
\ref{ch:25}장 가 이 격차를 \emph{데이터 축을 통제하고} 좁히는 장입니다 - \inlinecode{gpt2} (124M) 의 사전학습 \emph{위에} 같은 TinyStories 30K 로 \textbf{continual pretraining}. \emph{대규모 일반 사전학습 모델을 작은 도메인 데이터로 적응} 시킬 때의 generation 품질이, 우리 from-scratch 작은 GPT 와 어떻게 다른지 직접 비교.
\end{quote}



\section{변형 - sampling hyperparam 비교}

같은 prompt 에 \inlinecode{temperature\ /\ top\_k\ /\ top\_p} 만 바꿔 generation 스타일 변화 관찰. \emph{학습된 본체는 그대로} - 변하는 건 \emph{sampling 분포} 뿐.



\Needspace{11\baselineskip}
\begin{lstlisting}
prompt = "Once upon a time, a little rabbit"
configs = [
    {"label": "T=0.3, top_k=20  (conservative)", "temperature": 0.3, "top_k": 20,  "top_p": None},
    {"label": "T=0.8, top_k=50  (balanced)",    "temperature": 0.8, "top_k": 50,  "top_p": None},
    {"label": "T=1.0, top_p=0.9 (nucleus)",     "temperature": 1.0, "top_k": 0,   "top_p": 0.9},
    {"label": "T=1.2, top_k=100 (diverse)",     "temperature": 1.2, "top_k": 100, "top_p": None},
]
for c in configs:
    torch.manual_seed(SEED)
    print("=" * 70)
    print(f"[{c['label']}]")
    print(generate_text(model, prompt, max_new_tokens=60, do_sample=True,
                        temperature=c["temperature"], top_k=c["top_k"], top_p=c["top_p"]))
    print()
\end{lstlisting}

\noindent\textbf{출력.}
\begin{bookoutputbox}
======================================================================
[T=0.3, top_k=20  (conservative)]
Once upon a time, a little rabbit named Lily. She loved to play with her to...

Lily was sad

======================================================================
[T=0.8, top_k=50  (balanced)]
Once upon a time, a little rabbit who liked to play in the park. One day, h...

One day, Timmy saw a small house, Timmy's mom a shiny tree in a big tree. T...

======================================================================
[T=1.0, top_p=0.9 (nucleus)]
Once upon a time, a little rabbit who liked to read water to play the squir...

...
\end{bookoutputbox}
\par\vspace{0.9em}
\textbf{관전 포인트}

\begin{itemize}
\tightlist
\item
  \inlinecode{temperature} 증가 \(\to\) logits 분포 \emph{평탄화} \(\to\) 다양성 증가, 일관성 감소
\item
  \inlinecode{top\_k=20} \(\to\) 매 step 후보를 \emph{상위 20 개} 로만 한정 \(\to\) 안전하지만 반복적
\item
  \inlinecode{top\_p=0.9} (nucleus) \(\to\) 누적 확률 90\% 이내 후보 \(\to\) \emph{모델이 확신 있을 땐 좁게, 애매할 땐 넓게} 자동 조정
\item
  \inlinecode{T=1.2,\ top\_k=100} \(\to\) 가장 다양하지만 \emph{말이 안 되는 토큰} 도 종종 섞임
\end{itemize}



\section{이 장에 등장한 라이브러리·함수}

\begin{booktable}{24장 새로 등장한 라이브러리}{tab:ch24-13}
\begin{adjustbox}{max width=\textwidth}
\begin{tabular}{@{}lll@{}}
\toprule
이름 & 한 줄 설명 & 다음 장에서 \\
\midrule

\inlinecode{transformers.GPT2Config} & GPT-2 구조 hyperparam (n\_layer, n\_head, n\_embd, n\_positions, ...) & \ref{ch:25}장 - \inlinecode{gpt2} 사전학습 config 그대로 로드 \\
\inlinecode{transformers.GPT2LMHeadModel} & decoder + LM head 내장, causal attention 자동 처리 & \ref{ch:25}장 - \inlinecode{AutoModelForCausalLM.from\_pretrained("gpt2")} \\
\inlinecode{transformers.GPT2LMHeadModel(config)} (random init) & from scratch 사전학습 모델 생성 & \ref{ch:26}장 - 한국어 TinyStories scratch \\
\inlinecode{tokenizers.BPE\ +\ ByteLevel} & byte-level BPE 토크나이저 직접 학습 & \ref{ch:26}장 - 한국어 BBPE \\
\inlinecode{DataCollatorForLanguageModeling(mlm=False)} & CausalLM collator - \inlinecode{labels\ =\ input\_ids.clone()} 자동 & \ref{ch:25}--\ref{ch:26}장 동일 / \ref{ch:28}장 SFT 는 \emph{\inlinecode{-100} 자리만 다름} \\
\inlinecode{group\_texts} 패턴 (HF run\_clm.py 표준) & 가변 길이 텍스트 \(\to\) 고정 length 토큰 블록 스트림 & \ref{ch:25}--\ref{ch:26}장 동일 \\
\inlinecode{model.generate(do\_sample=True,\ ...)} & sampling-based text generation (temperature / top\_k / top\_p) & \ref{ch:25}--\ref{ch:30}장 전반에서 활용 \\
\texttt{\textless{}\textbar{}endoftext\textbar{}\textgreater{}} 단일 special token & GPT-2 컨벤션 (bos = eos = pad 겸용) & \ref{ch:25}장 - 같은 컨벤션 \\
\bottomrule
\end{tabular}
\end{adjustbox}
\end{booktable}



\section{체크포인트 질문}

\begin{enumerate}
\def\labelenumi{\arabic{enumi}.}
\tightlist
\item
  GPT CausalLM 사전학습은 \emph{거의 모든 자리} 가 학습 신호인데, BERT MLM 은 \emph{15\% 자리} 만 학습 신호입니다. 왜 BERT 는 \emph{전 자리} 를 학습하지 못할까요? (causal attention vs bidirectional attention 의 정보 누출 관점)
\item
  \inlinecode{tie\_word\_embeddings=True} (weight tying) 가 \inlinecode{Linear(H,\ V)} 의 파라미터를 어떻게 절약하나요? vocab 2,048 / hidden 256 일 때 절약되는 파라미터 수를 직접 계산해 보세요.
\item
  학습 첫 step loss 가 \(\ln(2048) \approx 7.62\) 가 아니라 \emph{5.0} 이라면 무엇을 의심해야 하나요? \emph{15.0} 이라면?
\item
  \ref{ch:28}장 (SFT) 에서는 \texttt{labels{[}:prompt\_len{]}\ =\ -100} 한 줄로 \emph{prompt 부분만 학습 신호에서 제외} 합니다. 이 장의 collator 출력 (거의 모든 자리가 학습 신호) 과 비교해 \emph{왜 이 한 줄이 모델이 instruction 을 따라가게 만드는지} 설명해 보세요.
\end{enumerate}



\section{FAQ}

\begin{faqBox}{Q1. (이론) 왜 GPT 는 \emph{거의 모든 토큰} 을 학습 신호로 쓰고 BERT 는 15\% 만 쓰나요?}

\textbf{causal attention vs bidirectional attention 의 정보 누출 차이} 때문입니다.

BERT 의 \emph{bidirectional} attention 은 토큰 \(i\) 의 hidden 이 좌·우 모든 토큰을 다 봅니다. 만약 \emph{모든 자리} 의 정답 토큰을 \emph{예측 task} 로 두면, 모델은 \emph{자기 자신을 그대로 복사} 하는 trivial 해를 학습합니다 (input 이 hidden 에 그대로 들어 있으니까). 그래서 BERT 는 일부 토큰을 \texttt{{[}MASK{]}} 로 \emph{가려야만} 의미 있는 학습 신호가 생깁니다 - \emph{주변 문맥} 으로 \emph{가려진 자리} 를 복원.

GPT 의 \emph{causal} attention 은 토큰 \(i\) 의 hidden 이 \emph{과거 (j \(\le\) i)} 만 봅니다. 미래 토큰을 못 보니 \emph{next-token 예측} 이 trivial 하지 않습니다 - 모든 자리에서 \emph{다음 토큰} 을 예측해도 cheating 이 안 됩니다. 그래서 \emph{전 자리} 가 학습 신호.

코드 한 줄로 갈리는 차이:

\Needspace{11\baselineskip}
\begin{lstlisting}
# BERT MLM
DataCollatorForLanguageModeling(
    tokenizer,
    mlm=True,
    mlm_probability=0.15,
)

# GPT CausalLM
DataCollatorForLanguageModeling(tokenizer, mlm=False)
\end{lstlisting}

한 step 의 \emph{토큰 학습 효율} 은 GPT 가 약 5-6배 높습니다 (15\% vs 거의 100\%). 그래서 같은 step 수라도 GPT 가 더 많은 토큰을 학습.

\end{faqBox}

\begin{faqBox}{Q2. (이론) TinyStories 는 \emph{일반 도메인} 인가요 \emph{task corpus} 인가요? 왜 일반 위키 (Wikitext-103) 가 아닌가요?}

\textbf{TinyStories 는 \emph{합성된 simple 스토리} 라 \emph{generation 시연 가치} 가 우선} 인 데이터입니다. \emph{진정한 일반 도메인 사전학습} 의 의미에서는 Wikitext-103 보다 약하지만, 본 장의 목적은 \emph{작은 모델로 generation 이 어떻게 동작하는지를 직접 보는 것} - 일반 위키 (Wikitext-103) 로 같은 셋업을 돌리면 3M 모델이 \emph{문장 구조를 학습하기 전에 학습이 끝남}. TinyStories 의 단순한 어휘·문법 덕분에 \emph{작은 모델로도 grammatical 한 생성이 가능} 합니다.

\ref{ch:25}장 가 그 \emph{trade-off 의 반대편} 을 다룹니다 - \emph{큰 모델 (gpt2 124M) + 대규모 일반 코퍼스 (WebText)} 의 사전학습된 본체를 TinyStories 로 \textbf{continual pretraining}. \emph{작은 + 합성 도메인 from-scratch} vs \emph{큰 + 일반 도메인 사전학습 후 continual pretraining} 의 generation 품질 격차가 핵심 비교.

\end{faqBox}

\begin{faqBox}{Q3. (이론) BPE 토크나이저는 \ref{ch:19}장의 WordPiece / WordLevel 과 어떻게 다른가요?}

세 방식 모두 \emph{vocab 학습 알고리즘} 이지만, \emph{어떻게 subword 를 만드는가} 가 다릅니다.

\begin{itemize}
\tightlist
\item
  \textbf{WordLevel}: 공백 + 빈도 - 단어 통째로 vocab 등록. UNK 가 많음.
\item
  \textbf{WordPiece (BERT)}: 빈도 + likelihood 기반 subword 병합. 단어 \emph{중간} subword 에 \inlinecode{\#\#} 접두사 (\inlinecode{unhappiness} \(\to\) \texttt{{[}“un”,\ “\#\#happiness”{]}}).
\item
  \textbf{BPE (GPT-2)}: \emph{byte 쌍 빈도} 기반 반복 병합. 접두사 없이 byte 시퀀스 그대로 (\inlinecode{unhappiness} \(\to\) \texttt{{[}“un”,\ “happiness”{]}}).
\end{itemize}

본 장은 \emph{byte-level BPE} - 가장 작은 단위가 \emph{byte (256개)} 라 \emph{어떤 유니코드 문자열} (이모지, 한글, 특수 기호) 도 UNK 없이 표현 가능. \ref{ch:19}장의 직접 학습 장에서 본 세 알고리즘 중 \emph{GPT 계열이 BPE 를 선호} 하는 이유.

\Needspace{11\baselineskip}
\begin{lstlisting}
from tokenizers.models import BPE
from tokenizers.trainers import BpeTrainer
from tokenizers.pre_tokenizers import ByteLevel

bpe = Tokenizer(BPE(unk_token=None))
bpe.pre_tokenizer = ByteLevel(add_prefix_space=False)
trainer = BpeTrainer(
    vocab_size=2048,
    initial_alphabet=ByteLevel.alphabet(),
)
bpe.train_from_iterator(text_iter, trainer)
\end{lstlisting}

\end{faqBox}

\begin{faqBox}{Q4. (이론) \inlinecode{tie\_word\_embeddings=True} (weight tying) 가 정확히 무엇을 공유하나요?}

\textbf{Input embedding (\inlinecode{wte}) 의 weight 와 LM head (\inlinecode{Linear(H,\ V)}) 의 weight 가 \emph{같은 텐서를 공유}} 합니다.

\Needspace{5\baselineskip}
\begin{lstlisting}
model.lm_head.weight = model.transformer.wte.weight
\end{lstlisting}

직관: input embedding 은 \emph{vocab token \(\to\) hidden} 변환, LM head 는 \emph{hidden \(\to\) vocab logit} 변환. 둘이 \emph{transpose 관계} 라 같은 weight 를 공유해도 의미가 통합니다. 효과:

\begin{itemize}
\tightlist
\item
  \textbf{파라미터 절약}: \(\texttt{vocab\_size} \times \texttt{hidden\_size}\) 만큼입니다. 본 장에서는 \(2{,}048 \times 256 = 524{,}288\), 즉 약 0.5M params 입니다. 전체 3M 모델의 약 17\%라 작은 모델에서는 비중이 큽니다.
\item
  \textbf{학습 안정}: input 과 output 이 \emph{같은 임베딩 공간} 을 공유 \(\to\) 일관성이 높아집니다.
\end{itemize}

GPT-2 의 기본값. 우리 모델도 자동으로 적용됩니다 (\inlinecode{config.tie\_word\_embeddings=True}).

\end{faqBox}

\begin{faqBox}{Q5. (실무) temperature / top\_k / top\_p sampling 의 의미는?}

\inlinecode{model.generate(do\_sample=True,\ ...)} 의 세 핵심 hyperparam:

\begin{itemize}
\tightlist
\item
  \textbf{\inlinecode{temperature}} (T): logits 를 \emph{나눠} softmax \(\to\) T\textless1 은 분포를 뾰족하게 (안전), T\textgreater1 은 평탄화 (다양). \(p_i = \text{softmax}(\text{logits}_i / T)\).
\item
  \textbf{\inlinecode{top\_k}}: 매 step \emph{상위 k 개} 후보로만 한정. 작으면 안전하지만 반복적.
\item
  \textbf{\inlinecode{top\_p}} (nucleus): \emph{누적 확률 p} 이내 후보로 한정. 모델이 확신 있을 땐 좁게 (top-1 이 0.9), 애매할 땐 넓게 (top-100 이 0.9) 자동 조정.
\end{itemize}

\Needspace{11\baselineskip}
\begin{lstlisting}
# 1. logits / T  $\to$ softmax

model.generate(
    do_sample=True,
    temperature=0.8,
    top_k=50,
    top_p=0.9,
    max_new_tokens=60,
)
\end{lstlisting}

일반적 추천: \emph{대화} 에는 \inlinecode{T=0.7-0.9,\ top\_p=0.9}, \emph{코드 / 수식} 에는 \inlinecode{T=0.2-0.4,\ top\_k=20}, \emph{창의적 생성} 에는 \inlinecode{T=1.0-1.2,\ top\_p=0.95}.

\end{faqBox}

\begin{faqBox}{Q6. (실무) \inlinecode{labels\ =\ -100} 트릭이 CausalLM 사전학습에서 거의 안 쓰이는데, 그럼 언제 쓰이나요?}

본 장의 collator 출력에서 봤듯 CausalLM 사전학습은 \emph{pad 토큰 자리만} \inlinecode{-100} (그것도 \inlinecode{group\_texts} 로 chunk 길이가 모두 같으면 pad 도 없음). 거의 안 쓰임.

하지만 같은 트릭이 \textbf{\ref{ch:28}장 (SFT, Instruction Tuning)} 에서 \emph{결정적 한 줄} 로 부활합니다:

\Needspace{9\baselineskip}
\begin{lstlisting}
prompt = "### 질문: 한국의 수도는?\n### 답변: "
response = "서울입니다."

input_ids = tokenizer(prompt + response)["input_ids"]
labels = input_ids.copy()
prompt_len = len(tokenizer(prompt)["input_ids"])
labels[:prompt_len] = [-100] * prompt_len
\end{lstlisting}

이 한 줄로 모델은 \emph{prompt 토큰을 외우지 않고 response 만 학습} - 같은 instruction 에 대한 \emph{다양한 response} 가 학습 가능하고, 추론 시에는 \emph{주어진 instruction 에 대해 response 를 생성} 하게 됩니다. \emph{모델이 instruction 을 따라간다} 는 게 이 한 줄의 효과.

본 장의 collator 출력 (거의 모든 자리 = 학습 신호) 을 손에 익혀 두면 \ref{ch:28}장의 \texttt{labels{[}:prompt\_len{]}\ =\ -100} 가 단번에 이해됩니다.

\end{faqBox}

\begin{faqBox}{Q7. (실무) 다음 장 (\ref{ch:25}장 --- continual pretraining) 와의 비교는 어떻게 되나요?}

\ref{ch:25}장 = \emph{OpenAI 가 사전학습한 \inlinecode{gpt2} (124M params, WebText 약 40GB) 를 TinyStories 로} \textbf{continual pretraining} (같은 CausalLM task, 같은 LM head --- \emph{task adaptation 의미의 fine-tune 이 아님}). 본 장의 \emph{작은 from-scratch 모델} 과 \emph{완전 반대 출발점}:

\begin{booktable}{24장 FAQ 표}{tab:ch24-14}
\begin{adjustbox}{max width=\textwidth}
\begin{tabular}{@{}lll@{}}
\toprule
축 & \ref{ch:24}장 (본 장) & \ref{ch:25}장 (다음) \\
\midrule

모델 크기 & 약 3M params & \textbf{약 124M (40배)} \\
사전학습 & from scratch (random init) & \textbf{OpenAI WebText 약 40GB 사전학습} \\
TinyStories 학습 & 사전학습 그 자체 (1500 steps) & \textbf{continual pretraining} (약 3,000 step 규모) \\
토크나이저 & 직접 학습 BPE vocab 2,048 & \textbf{gpt2 BPE vocab 50,257 (그대로)} \\
Generation 품질 & grammatical 한 동화 풍 영어 & \textbf{자연스러운 동화 + 일반 도메인 폭} \\
학습 시간 & 약 1분 (사전학습) & \textbf{약 19분} (continual pretraining)  \\
\bottomrule
\end{tabular}
\end{adjustbox}
\end{booktable}

\textbf{핵심 메시지}: \emph{대규모 일반 사전학습된 본체} + \emph{작은 도메인 continual pretraining} 이 \emph{작은 from-scratch 모델} 보다 \emph{빠르게, 그리고 좋게} 도달합니다. \emph{왜 실무는 보통 from-scratch 가 아니라 사전학습 모델을 가져와 새 데이터로 계속 학습하거나 SFT 하는가} 의 답. (단계 3 SFT 는 \ref{ch:28}장에서 본격.)
\end{faqBox}



\begin{previewBox}{미리보기: 다음 장}

\textbf{\ref{ch:25}장. GPT2 (124M) Continual Pretraining 으로 TinyStories 에 적응 --- \emph{대규모 사전학습 모델의 도메인 계속 학습}}

\begin{itemize}
\tightlist
\item
  \inlinecode{AutoModelForCausalLM.from\_pretrained("gpt2")} - OpenAI WebText 약 40GB 로 사전학습된 124M params 모델 로드
\item
  \textbf{같은 TinyStories 30K} 데이터 (본 장와 동일) 로 \textbf{continual pretraining} (계속 사전학습 / continual learning --- \emph{같은 CausalLM task, 새 데이터, head 그대로}. \emph{task adaptation 의미의 fine-tune 이 아니라 단계 2})
\item
  \textbf{핵심 비교}: 본 장 (3M, from scratch, T4 약 1분) vs \ref{ch:25}장 (124M, continual pretraining, T4 약 19분) 의 generation 품질·학습 곡선 격차
\item
  \emph{trainer 자체는 \ref{ch:24}장과 동일} (\inlinecode{transformers.Trainer} + \inlinecode{DataCollatorForLanguageModeling(mlm=False)}) --- \emph{변하는 건 모델 로드 한 줄 + lr (scratch 5e-4 \(\to\) continual pretraining 2e-5)}
\item
  작은 데이터 + 큰 사전학습 모델 = \emph{왜 실무가 from-scratch 가 아니라 사전학습 모델 위에 계속 학습 패턴인가} 의 정량 답변
\item
  \emph{진짜 task adaptation 의미의 fine-tune (instruction tuning)} 은 \ref{ch:28}장 SFT 에서 본격 등장
\end{itemize}

\begin{quote}
\textbf{변하는 축}: \emph{모델 크기 + 사전학습 여부} (3M / scratch \(\to\) 124M / pretrained). 데이터·토크나이저 규약·loss·trainer 는 같음. Phase 4 의 \emph{학습 단계 2 (continual pretraining)} 가 본격적으로 자리 잡는 장.
\end{quote}
\end{previewBox}
